% =============================================================================
%  Bubble Radius Fitting — Quick Start Guide
%  Author: Zixiang (Zach) Tong, The University of Texas at Austin
% =============================================================================
\documentclass[11pt, letterpaper]{article}

\usepackage[margin=1in]{geometry}
\usepackage{enumitem}
\usepackage{xcolor}
\usepackage{hyperref}
\usepackage{fancyhdr}
\usepackage{titlesec}
\usepackage{booktabs}
\usepackage{graphicx}
\usepackage{tcolorbox}

% ---------- Colours ----------
\definecolor{primaryblue}{HTML}{3B82F6}
\definecolor{darknavy}{HTML}{0F172A}
\definecolor{successgreen}{HTML}{22C55E}
\definecolor{errorred}{HTML}{EF4444}
\definecolor{dimgrey}{HTML}{94A3B8}
\definecolor{lightslate}{HTML}{F8FAFC}

\hypersetup{
    colorlinks=true,
    linkcolor=primaryblue,
    urlcolor=primaryblue,
}

\pagestyle{fancy}
\fancyhf{}
\fancyhead[L]{\small\textcolor{dimgrey}{Bubble Radius Fitting --- Quick Start Guide}}
\fancyhead[R]{\small\textcolor{dimgrey}{\thepage}}
\renewcommand{\headrulewidth}{0.4pt}

\titleformat{\section}
  {\Large\bfseries\color{darknavy}}{\thesection}{1em}{}
\titleformat{\subsection}
  {\large\bfseries\color{darknavy}}{\thesubsection}{1em}{}

\newtcolorbox{tipbox}[1][]{
    colback=blue!5, colframe=primaryblue,
    fonttitle=\bfseries, title={#1},
    boxrule=0.5pt, arc=4pt,
}

% =============================================================================
\begin{document}

% -------- Title --------
\begin{center}
    {\Huge\bfseries\color{darknavy} Bubble Radius Fitting}\\[0.4cm]
    {\LARGE\color{primaryblue} Quick Start Guide}\\[0.8cm]
    {\large Zixiang (Zach) Tong\quad---\quad The University of Texas at Austin}\\[0.3cm]
    {\large \today}
\end{center}

\vspace{0.5cm}
\hrule
\vspace{0.8cm}


% =============================================================================
\section{Launch the Application}
% =============================================================================

\textbf{Option A --- From source:}
\begin{verbatim}
    cd bubbletrack
    .venv\Scripts\activate
    bubbletrack
\end{verbatim}

\textbf{Option B --- Standalone:} Double-click \texttt{BubbleRadiusFitting.exe}.


% =============================================================================
\section{Load Images}
% =============================================================================

\begin{enumerate}
    \item In the left sidebar, click \textbf{Browse}.
    \item Select the folder containing your image sequence
          (\texttt{.tiff/.tif/.png/.jpg/.bmp}).
    \item The first image loads automatically. Use the \textbf{frame scrubber}
          slider at the bottom to browse frames.
\end{enumerate}


% =============================================================================
\section{Set Up Parameters (Pre-tune Tab)}
% =============================================================================

\begin{enumerate}
    \item \textbf{Select ROI:} Click \textbf{Select ROI}, then drag a
          rectangle on the image to enclose the bubble.

    \item \textbf{Adjust Threshold:} Drag the slider until the bubble
          interior is clearly white in the binary mask panel on the right.

    \item \textbf{Adjust Removing Factor:} Increase to remove small noise
          spots from the binary image.

    \item \textbf{Edge Flags:} If the bubble extends beyond the ROI boundary,
          check the corresponding edge (Top / Right / Down / Left).

    \item \textbf{Advanced Filters} (optional, expand the collapsible
          section):
          \begin{itemize}
              \item \textbf{Gaussian Blur} --- smooths noise before
                    binarisation (e.g.\ tracer particles).
              \item \textbf{CLAHE} --- improves contrast when lighting is
                    uneven.
              \item \textbf{Morph Close} --- fills small holes inside the
                    bubble boundary.
              \item \textbf{Morph Open} --- removes thin spurs on the bubble
                    edge.
          \end{itemize}

    \item \textbf{Test:} Click \textbf{Fit Current Frame}. A
          \textcolor{primaryblue}{blue circle} and
          \textcolor{errorred}{red edge points} should appear on the original
          image.
\end{enumerate}

\begin{tipbox}[Tip]
    Click \textbf{Save Preset} to store your tuned parameters to a JSON file.
    Use \textbf{Load Preset} to restore them in future sessions.
\end{tipbox}


% =============================================================================
\section{Batch Processing (Automatic Tab)}
% =============================================================================

\begin{enumerate}
    \item Switch to the \textbf{Automatic} tab.
    \item Set the \textbf{From} and \textbf{To} frame numbers (default: all
          frames).
    \item Click \textbf{Fit All}. The progress bar, image panels, and $R(t)$
          chart update in real time.
    \item Click \textbf{Stop} at any time to pause (results so far are kept).
    \item Click \textbf{Clear} to discard all results and start over.
\end{enumerate}


% =============================================================================
\section{Manual Fitting (Manual Tab)}
% =============================================================================

Use this when automatic detection does not work well for a particular frame.

\begin{enumerate}
    \item Switch to the \textbf{Manual} tab.
    \item Click \textbf{Select Points}, then click 3 or more points on the
          bubble edge (5--10 recommended).
    \item Click \textbf{Done} to fit a circle through the selected points.
    \item Click \textbf{Clear} to reset and try again.
\end{enumerate}


% =============================================================================
\section{Export Results}
% =============================================================================

In the \textbf{Post Processing} section at the bottom of the sidebar:

\begin{enumerate}
    \item Set \textbf{FPS} (frame rate) and its unit (Hz / kHz / MHz).
    \item Set \textbf{Scale} (\textmu m/px).
    \item Set \textbf{Rmax Fit Length} (odd number; window for peak
          interpolation).
\end{enumerate}

Then click one of the two export buttons:

\begin{center}
\renewcommand{\arraystretch}{1.4}
\begin{tabular}{lp{8cm}}
    \toprule
    \textbf{Button} & \textbf{What it saves} \\
    \midrule
    \textbf{Export Pixel Data} &
        Radius (px), circle centres, and edge points for every frame.
        Default filename: \texttt{radius\_pixel.mat} \\
    \textbf{Export Physical Data} &
        Time (s) vs.\ Radius (\textmu m) with $R_{\max}$ quadratic
        interpolation.
        Default filename: \texttt{radius\_time\_physical.mat} \\
    \bottomrule
\end{tabular}
\end{center}

A save dialog will open; choose the location and filename, then confirm.


% =============================================================================
\section{Loading Results in MATLAB}
% =============================================================================

\begin{verbatim}
    % Pixel data
    data = load('radius_pixel.mat');
    R_px = data.Radius;           % 1 x N, pixels (-1 = unfitted)

    % Physical data
    data = load('radius_time_physical.mat');
    plot(data.t * 1e6, data.R, 'r.');
    xlabel('Time (\mus)');
    ylabel('Radius (\mum)');
\end{verbatim}


% =============================================================================
\section{Common Issues}
% =============================================================================

\begin{center}
\renewcommand{\arraystretch}{1.4}
\begin{tabular}{p{5cm}p{7.5cm}}
    \toprule
    \textbf{Problem} & \textbf{Fix} \\
    \midrule
    Binary mask all white / all black &
        Adjust the Threshold slider. \\
    Noise spots in binary mask &
        Increase Removing Factor; enable Gaussian Blur. \\
    ``Too few edge points'' &
        Enlarge ROI; adjust Threshold; check edge flags. \\
    ``Radius outlier, skipped'' &
        Tune Threshold or filters to get a cleaner boundary. \\
    Gaps in bubble boundary &
        Enable CLAHE and/or Morph Close. \\
    Export error about fit length &
        Reduce Rmax Fit Length or process more frames. \\
    \bottomrule
\end{tabular}
\end{center}

\vfill
\begin{center}
    \small\textcolor{dimgrey}{%
        Bubble Radius Fitting --- developed at The University of Texas at Austin
    }
\end{center}

\end{document}
